\chapter{Noise, crosstalk and molecular fault models}

Chapter 22 ended with a deliberately uncomfortable result: a tile system can be
directed in the abstract model and still produce defective molecular objects. That
gap is not a footnote. It is the engineering problem.

Reliability begins by naming the event being counted. A \Term{growth error} is not
the same event as \Term{spurious nucleation}; a malformed origami is not necessarily
a wrong logical bit; an AFM classification error is not a molecular attachment
error. If these events share one number called ``the error rate'', neither a model
nor an experiment can tell us what improved.

\section{A reliability stack}

Let $C$ be the intended computational construction, $A$ the realized molecular
assembly, $Y$ the instrument observation and $\hat D$ a classified defect label.
The experimental pipeline is therefore
\[
 C \longrightarrow A \longrightarrow Y \longrightarrow \hat D.
\]
The first mapping is conditioned on design $C$: $P(A\mid C)$, not a promise of
physical realization. Measurement has channel $P(Y\mid A)$ and classification
has $P(\hat D\mid Y)$. Composition gives
\[
P(\hat D\mid C)=\sum_A\int P(\hat D\mid Y)\,P(Y\mid A)\,P(A\mid C)\,dY.
\]
Use a sum rather than the integral for a discrete observation. This factorization
assumes the declared intermediates include the relevant history; an instrument that
depends on handling condition needs that variable too. Formal correctness constrains
the intended object. Chemistry determines the distribution over $A$. Measurement maps
physical objects to signals. Analysis maps signals to labels.

\Plate{01-stack}{Reliability is an end-to-end contract. A defect can be introduced
by assembly or by observation, so the experiment must define both molecular and
measurement events.}{dnad23:stack}

Suppose an assembly contains $N$ positions at which a particular local error could
occur. Define $E_i=1$ if site $i$ is erroneous under a declared classifier.
The local error fraction is
\[
\hat p={1\over N}\sum_{i=1}^N E_i.
\]
This is an estimator with a denominator. It is not automatically the probability
that an entire object is unusable. A single local error may be benign, catastrophic,
or invisible to the measured function.

\section{Why size amplifies small local error}

As a first teaching model, assume independent identical site errors with probability
$p$. Then
\begin{equation}
P(\hbox{all $N$ sites correct})=(1-p)^N.
\label{dnad23:iid}
\end{equation}
At $p=0.01$, one hundred sites give about $0.366$ all-correct probability; ten
thousand sites drive it essentially to zero. These values are generated by
\texttt{reference.py}.

\Plate{02-yield}{The IID law explains finite-size amplification but is not a
universal yield law. Correlated defects, growth-front history and reversible
correction can change the distribution.}{dnad23:yield}

The assumption matters as much as the equation. If each object first draws a latent
error propensity $q$ and its sites then fail conditionally independently with
probability $q$, errors are positively correlated within objects. A beta distribution
for $q$ gives a beta-binomial count model. The package compares an IID model and a
beta-binomial model with the same mean local error and obtains different all-correct
probabilities. Thus ``we measured one percent errors'' does not determine object
yield without dependence information.

\subsection{Same marginals, different products}
Set $K=\sum_i E_i$, and suppose $q\sim\operatorname{Beta}(\alpha,\beta)$ with
$\alpha,\beta>0$. Conditional on $q$, $K\sim\operatorname{Binomial}(N,q)$.
Integrating the binomial polynomial against the beta density gives
\begin{equation}
P(K=k)={N\choose k}
\frac{B(k+\alpha,N-k+\beta)}{B(\alpha,\beta)},\qquad 0\le k\le N.
\label{dnad23:beta}
\end{equation}
Each site has marginal $p=\alpha/(\alpha+\beta)$. For two distinct sites,
$\operatorname{Cov}(E_i,E_j)=\operatorname{Var}(q)$, so their correlation is
$\rho=1/(\alpha+\beta+1)$. Therefore
\[
\operatorname{Var}(K)=Np(1-p)\,[1+(N-1)\rho].
\]
Treating clustered sites as independent underestimates this model's count variance.
For all-correct probability,
\[
P(K=0)=\mathbb E[(1-q)^N]
=\prod_{i=0}^{N-1}\frac{\beta+i}{\alpha+\beta+i}.
\]
For $N\ge2$, $(1-q)^N$ is convex. Jensen's inequality implies this mixture has
at least the IID all-correct probability at its mean $p$. This does \emph{not}
mean correlation is beneficial: it creates more very clean objects and more
concentrated damage. The variance and the logical task both matter.

\Listing{all_correct_beta_binomial}
\begin{center}\input{\Generated/yield_examples.tex}\end{center}
The table fixes $p=0.01$, $\alpha=0.99$, $\beta=98.01$ and varies only object
size. Values are computed model probabilities. No DNA experiment supplied these
parameters. With no dependence assumption at all, the union bound and containment
in each site's correct event give
\[
\max(0,1-Np)\le P(K=0)\le1-p \quad (N\ge1).
\]
The empty object has all-correct probability one. These wide bounds expose how
much an IID assertion contributes beyond a measured marginal fraction.

\section{Redundancy as a mathematical mechanism}

The simplest redundancy model encodes one logical bit in an odd number $r$ of
independent physical bits and decodes by majority. If each physical bit flips with
probability $p$, the exact logical failure probability is
\begin{equation}
P_r(p)=\sum_{k=(r+1)/2}^{r}{r\choose k}p^k(1-p)^{r-k}.
\label{dnad23:majority}
\end{equation}
For $r=3$,
\[
P_3(p)=3p^2(1-p)+p^3=3p^2-2p^3.
\]
For small $p$, the leading term is $3p^2$: one error is correctable and two are
required for failure. The test suite checks this closed form and compares the exact
five-copy result with an independently sampled Monte Carlo oracle.

\Listing{majority_failure}
The odd-copy decoder has no tie convention. At $p=1/2$ its error is $1/2$;
above $1/2$ redundancy amplifies the wrong majority. Its apparent benefit below
$1/2$ also depends on the independence of the copies.
\begin{center}\input{\Generated/repetition.tex}\end{center}
In the correlated column, all copies share the beta-distributed propensity
above. The logical tail is the sum of Equation~\ref{dnad23:beta} over
$k>(r/2)$, not the binomial tail evaluated at mean $p$. In a more extreme
common-mode counterexample, all $r$ copies flip together with probability $p$.
Majority failure is then exactly $p$ for every odd $r$. More material has not
created more independent evidence. This is why physically colocated redundancy
needs a model of its shared environment, not only a copy count.

\Plate{03-repetition}{Majority redundancy converts several physical observations
into one logical decision. This mathematical decoder illustrates error suppression;
molecular proofreading must realize redundancy through assembly pathways rather
than by assuming an external perfect voter.}{dnad23:repeat}

This distinction prevents a common conceptual shortcut. A DNA tile crystal does not
contain a free digital majority gate merely because our notebook can decode three
pixels. Molecular redundancy must alter the physical pathway so that a wrong logical
choice requires multiple coordinated unfavorable events.

\section{Proofreading tile sets}

Proofreading constructions replace a logical tile by a structured block of physical
tile types. Internal glues are chosen so that completing an incorrect logical block
requires additional unfavorable attachments; correct growth can exploit cooperative
support. Winfree and Bekbolatov analyzed such constructions with kinetic and
thermodynamic tile models and predicted substantially improved error scaling over
direct implementations~\cite{WinfreeBekbolatov2004}.

\Plate{04-proofreading}{Proofreading expands a logical tile into a cooperative
physical structure. The cost is explicit: more tile identities/instances, larger
area and altered growth kinetics.}{dnad23:proof}

A useful abstraction is an \Term{error exponent}. If a direct construction has
logical error $O(\epsilon)$ while a redundant construction requires two effectively
independent rare events, its leading error can become $O(\epsilon^2)$. The exponent
is a mechanistic statement only under the model's hypotheses. Correlated wrong
attachments, facet nucleation or kinetic trapping can restore lower-order failure
paths.

Experiments matter because they test whether the pathway survives chemistry.
Schulman, Wright and Winfree later built DNA double-crossover systems in which a
short bit string was redundantly encoded by one through four molecules. Under their
reported conditions, each added redundancy level reduced bitwise error by roughly a
factor of three, and fourfold redundancy produced an error rate below 0.1 percent.
That is evidence for a particular physical design and condition—not a universal
``threefold per tile'' law~\cite{Schulman2015}.

\section{Reliability has several failure channels}

A \Term{mismatch} is an unintended local adjacency or logical identity at the growth
front. \Term{Facet nucleation} begins a new growth patch on a crystal facet where the
intended algorithm does not authorize one. \Term{Spurious nucleation} more broadly
means seed-independent initiation. A \Term{kinetic trap} is an off-pathway state
whose escape is slow relative to the experimental observation window.

\Plate{05-errors}{Different failure channels branch from different physical events.
Combining them into one undifferentiated rate can hide the mechanism that a
redundancy scheme actually changes.}{dnad23:errors}

Snaked proofreading was designed partly to make facet nucleation require a more
difficult sequence of attachments. In the reported experiment, its facet nucleation
rate was fourfold smaller than uniform proofreading under the tested conditions.
This result is especially instructive because the measured event is not the same
event as a bitwise growth mismatch~\cite{Chen2007}.

\section{Speed, stability and correction}

Error suppression is not free. Lower free-tile concentration can reduce some
undesired association events but slows desired collision rates. Stronger binding can
stabilize a correct frontier but can also prolong an incorrect trap. Greater
redundancy increases material, area and often dependency depth.

Let $k_{\rm on}c$ denote a schematic association hazard and $k_{\rm off}$ a
dissociation hazard for one context. The ratio of these hazards influences whether a
weakly supported attachment tends to disappear before it is locked in by subsequent
growth. But neither symbol is a universal constant: sequence, temperature, salt,
geometry and neighboring contacts matter. A reliability model should therefore
report the condition under which its rates were estimated.

There is also an observation-time tradeoff. A product scored too early may count
recoverable intermediates as defects; scored too late, it may accumulate degradation
or aggregation. ``Yield at time $t$'' is more reproducible than ``yield'' without a
time coordinate.

\section{Origami reliability and off-target interactions}

Origami adds failure modes beyond square-tile growth: missing staples, competing
partial folds, scaffold aggregation, topological frustration and off-target
hybridization. Recent experimental work has reinforced that the primary scaffold
sequence can alter folding reliability even when the intended complementarity and
target geometry are nominally preserved; excessive off-target interactions can
produce kinetic traps~\cite{ShirtEdiss2026}.

This is a direct lesson about abstraction. Two designs can have the same graph of
intended staple-scaffold contacts and still have different physical landscapes
because unintended sequence interactions differ. A design compiler that verifies
only intended complementarity is therefore incomplete as a reliability tool.

\section{Faults upstream of the assembly frontier}
Tile errors are only one instance of a broader molecular fault taxonomy.
A synthesis substitution changes an identity; an insertion or deletion changes
register and length; a missing component changes available multiplicity.
An unintended recognition interaction is \Term{crosstalk}: a pathway connects
objects that the design intended to keep separate. \Term{Leak} is output generated
without the intended input or authorization event. These definitions describe
different failures, even if each increases the same downstream fluorescence.

Separate a specification error from an implementation error. A wrong truth-table
entry is wrong even in a perfect abstract interpreter. An unintended sequence
contact can violate a correct symbolic specification. A stochastic correct-path
reaction can also miss its deadline without any wrong recognition. DNA strand
circuits, PCR-based workflows and tile assemblies consequently need different
fault models rather than one transferable scalar \emph{DNA noise}.

An elementary amplification counterexample illustrates composition bias without
pretending to simulate PCR chemistry. Let initial target fraction be $f\in[0,1]$,
and let two species multiply by constant per-cycle gains $g_T,g_B>0$. After $m$
cycles, ideal bookkeeping gives
\begin{equation}
f_m=\frac{f g_T^m}{f g_T^m+(1-f)g_B^m},\qquad
\frac{f_m}{1-f_m}=\frac{f}{1-f}\left(\frac{g_T}{g_B}\right)^m.
\label{dnad23:gain}
\end{equation}
The odds form applies to $0<f<1$; the endpoints remain zero and one.
Equal gains preserve composition, while a small ratio bias compounds across cycles.
The code evaluates log odds to avoid overflow. It assumes constant gains, no
resource depletion and no sequence-changing reactions; real amplification can
violate all three. Thus increased abundance is not evidence of unbiased recovery.
\Listing{amplification_fraction}
When a design is scored after amplification, distinguish assembly failure from
amplification or sampling bias. Chapter 25 develops storage/sequencing channels;
Chapter 26 supplies experiment records and controls.

\section{Measurement is part of the error model}

Suppose $n$ classified sites contain $k$ observed defects. Reporting $k/n$ without
uncertainty is weak evidence, especially when $k$ is small. The package computes a
Wilson interval for a binomial proportion as a transparent default:
\[
\tilde p={\hat p+z^2/(2n)\over1+z^2/n},\qquad
\Delta={z\sqrt{\hat p(1-\hat p)/n+z^2/(4n^2)}\over1+z^2/n}.
\]
If sites within one crystal are correlated, the crystal—not each pixel—may be the
appropriate independent replicate. Hierarchical or cluster-aware analysis is then
needed.

\subsection{Calibrate the classifier before inferring chemistry}
Let the true defect probability be $p$, the false-positive probability be
$\alpha_{\rm obs}=P(\hat D=1\mid D=0)$, and the false-negative probability be
$\beta_{\rm obs}=P(\hat D=0\mid D=1)$. These observation parameters are distinct
from the beta-mixture parameters. Total probability gives
\begin{equation}
p_{\rm obs}=\alpha_{\rm obs}
 +(1-\alpha_{\rm obs}-\beta_{\rm obs})p.
\label{dnad23:measurement}
\end{equation}
For $p=.01$, $\alpha_{\rm obs}=.02$, $\beta_{\rm obs}=.10$, the observed
defect rate is .0288. An assay can overstate a low physical error rate even
while missing some real defects.
\Listing{observed_defect_probability}
\Listing{infer_defect_probability}
If $1-\alpha_{\rm obs}-\beta_{\rm obs}=0$, observation carries no information
about $p$; inversion is not identifiable. Near zero, calibration errors are amplified.
An inverted value outside $[0,1]$ indicates inconsistent data or calibration,
not permission to silently clip a scientifically meaningful discrepancy.

\begin{center}\input{\Generated/wilson.tex}\end{center}
Zero observed defects in twenty independent trials still permits a sizeable upper
Wilson bound. This is an approximate interval for the \emph{observed} Bernoulli
event, not automatically for physical $p$. Calibrated uncertainty in both observation
parameters must also be propagated before making a chemical reliability claim.
\Listing{wilson}

Classification also needs controls. A blinded scorer should not know which design
is ``improved''. Positive controls establish that the assay can see a known defect;
negative controls estimate background classification. Multiple fields of view
should not be cherry-picked after inspection.

\Plate{06-loop}{A declared observation channel distinguishes false positives
from false negatives. Calibrated classification is needed before an observed
fraction can be interpreted as a physical defect probability.}{dnad23:loop}

\section{Reference implementation and independent checks}

The executable chapter checks three claims independently. First, exact binomial
majority failure is compared with Monte Carlo sampling. Second, the IID object-yield
law is contrasted with a correlated beta-binomial construction at the same mean
local error. Third, Wilson intervals expose the uncertainty of finite defect counts.
None of these calculations is labeled as a fitted DNA experiment.
Additional tests integrate the uniform-beta case analytically: its defect count is
uniform on $\{0,\ldots,N\}$, so each probability is $1/(N+1)$.
They check observation-channel inversion, the amplification odds identity, invalid
domains and extreme ratios. Monte Carlo agreement is a finite statistical check,
not an exact proof; its trial count and seed are recorded separately.

\section{A reliability ledger}

For each experiment, record at least:
\begin{enumerate}
\item the error event and denominator;
\item assembly condition, concentrations and observation time;
\item logical design and physical sequence design;
\item number of independently prepared samples and objects;
\item imaging/measurement protocol and classifier;
\item uncertainty and clustering assumptions;
\item speed, material and area costs of redundancy;
\item failure channels not measured by the chosen assay.
\end{enumerate}
This ledger makes comparisons harder to oversell and easier to reproduce.

\section{Exercises and worked solutions}
\begin{enumerate}
\item Derive Equation~\ref{dnad23:iid}. \textbf{Answer:} all $N$ independent
correct events each have probability $1-p$, so multiply them.
\item At $p=0.01$, derive $P_3$. \textbf{Answer:} failure requires exactly two or
three flips: $3p^2(1-p)+p^3$.
\item Why can correlated errors invalidate the IID yield law?
\textbf{Answer:} joint correctness no longer factorizes into marginal probabilities.
\item Is majority decoding a molecular proofreading mechanism?
\textbf{Answer:} not by itself; it is an external decoding rule unless the assembly
path physically realizes the redundancy and decision.
\item Name one cost of proofreading. \textbf{Answer:} more tile types/instances,
area, material, dependency depth or slower kinetics.
\item Distinguish mismatch and spurious nucleation. \textbf{Answer:} a mismatch is
a wrong local growth event on an assembly; spurious nucleation starts an unintended
assembly without the intended seed.
\item Why might stronger binding worsen reliability? \textbf{Answer:} it can make
incorrect trapped states slower to detach.
\item Why is ``0 defects in 20 objects'' not proof of zero error?
\textbf{Answer:} finite samples leave a nonzero upper confidence bound.
\item What is the independent unit if defects cluster within crystals?
\textbf{Answer:} often the crystal or preparation, requiring cluster-aware analysis.
\item How would you test whether off-target sequence interactions cause an origami
failure? \textbf{Answer:} hold target geometry/contact map as fixed as possible,
change sequence variants predicted to alter off-target interactions, use matched
conditions and blinded yield/defect measurements.
\item What would falsify a claimed proofreading benefit?
\textbf{Answer:} no reproducible reduction in the declared failure event after
matched conditions and uncertainty, or a benefit explained entirely by slower growth/
different concentration rather than the proposed cooperative mechanism.
\item Derive the beta-binomial correlation. \textbf{Answer:}
$\operatorname{Cov}(E_i,E_j)=\operatorname{Var}(q)$; divide by
$p(1-p)$ to obtain $1/(\alpha+\beta+1)$.
\item Give a redundancy failure with correct marginals. \textbf{Answer:} all
copies flip together with probability $p$; voting still fails with probability $p$.
\item Invert a .0288 observed defect rate at false-positive .02 and
false-negative .10. \textbf{Answer:} $(.0288-.02)/(.88)=.01$.
\item Why can no observed defects fail to imply no chemical defects?
\textbf{Answer:} finite independent samples leave uncertainty; a miscalibrated
classifier can also hide defects systematically.
\item What does equal amplification gain preserve? \textbf{Answer:} the
initial species fraction, under the two-species constant-gain assumptions.
\item Design a reliability benchmark. \textbf{Rubric:} predefine event and
denominator; include direct and redundant designs; matched physical conditions;
replicate preparations; time course; blinded classification; cluster-aware
uncertainty; report speed/material/area; separate growth, nucleation and observation
errors.
\end{enumerate}

\section{Bridge}
Reliability is not an extra layer added after molecular programming. It changes tile
sets, sequences, concentrations, timing, measurement and the meaning of a successful
run. Chapter 24 asks which resources scale with problem size, including material,
species, volume, latency and readout. Chapter 26 returns to experimental practice:
design records, calibration, controls, provenance and reproducible execution.
